\begin{EnAbstract}
% TEMPORARY DRAFT -- written from Chapters 1 to 3. Update after the evaluation (Chapter 4) is complete.
ECHONET Lite is an application-level communication protocol that is widely used in residential Home Energy Management Systems (HEMS) to connect controllers with appliances, smart meters and distributed energy resources. Its specification does not define message authentication, integrity protection or freshness validation within ECHONET Lite messages, and instead delegates security to lower layers on an optional and heterogeneous basis. As a result, a node that can reach the local network may construct, modify or replay messages that a receiving device accepts as legitimate, exposing household appliances to unauthorised control, replay and denial-of-service attacks. This thesis proposes the ECHONET Lite Authentication Framework (ELAF), an application-layer authentication mechanism that attaches a fixed-structure Security Extension Field (SEF) to each ECHONET Lite message without changing the protocol's object and property model. In its default configuration, the 23-byte SEF carries a version identifier, a length field, a mechanism identifier, a timestamp, a monotonic counter and an HMAC-SHA-256 authentication tag truncated to 96 bits. A receiving node accepts a message only if its structure, timestamp, counter and authentication tag are valid, and it maintains a request table that detects duplicated messages and suspends sources that exceed a configured request rate. The mechanism is implemented and evaluated in a containerised simulation environment and in an emulated testbed comprising commercial ECHONET Lite appliances, an ELAF-enabled controller and a translator node. The evaluation measures processing time, latency, throughput, detection accuracy, attack success rate and communication overhead under legitimate traffic and under injection, replay and flooding attacks.
% TODO(author): Add one to three sentences summarising the main quantitative results and the conclusion
% once Chapter 4 is complete. Do not add result values before they are measured.
\end{EnAbstract}
